\chapter{Scarcity is more than an empty store}\label{ch:scarcity}

\section{Water behind a locked gate}
A reservoir may contain enough water to supply a town today, yet its clinic can still have none. The connecting pipe may be broken, the water may be unsafe, the owner may deny access, or the clinic may be unable to pay the charge. Calling all these situations simply \emph{scarcity} hides the difference between them. A dry reservoir calls for a different response from a blocked pipe, and both differ from exclusion by price.

We can separate three questions. How much of a physical resource exists within the declared boundary? Can it be moved and used for the required function in time? Who is permitted and able to obtain it? A stock total answers the first question only. Infrastructure and quality enter the second; institutions and purchasing power enter the third. The FAO likewise treats water scarcity as involving physical supply, demand and access rather than one universal stock number.\cite{intro-water}

Some limits cannot be negotiated away. A town cannot distribute more drinkable water than it can obtain, treat or safely reuse over the relevant interval. Other limits are made by rules and may be changed. A full warehouse can stand beside hunger when the people who need food have no effective claim on it. This does not prove that all scarcity is artificial; it proves that physical abundance and social access are different measurements.

\section{Competition is a choice of allocation rule}
When several buyers want a constrained good, offers of money can decide who receives it. Competitive bidding is a clear way of rationing under a particular allocation rule. It can also make exclusion appear inevitable to those who cannot bid. Scarcity of the resource and scarcity of purchasing power then reinforce one another, although they are not the same thing.

The urge to gain an advantage is not a law of physics. Nor is cooperation an automatic consequence of shared need. People and organisations respond to the rules under which they operate. If an extra unit sold today increases private revenue while depletion tomorrow is borne by others, extracting and selling can be rewarded. If maintenance, restraint and reliable service are rewarded, the feasible choices look different. The design of incentives matters because it changes which physically possible actions are attractive.

This is why it would be mistaken to call maximisation itself the enemy. Every decision uses some criterion, even if it is not written down. The important question is \emph{what} is being maximised and \emph{subject to which conditions}. Maximising a short-term monetary return without a condition on essential services is not equivalent to maximising the continuing ability to meet needs. The first can grow while the second falls.

\section{Inequality changes exposure as well as income}
Unequal access is not merely a philosophical worry. The World Bank's \emph{Poverty, Prosperity, and Planet Report 2024} estimated that about 700 million people lived in extreme poverty under its then-used \$2.15-a-day line, and that progress in poverty reduction had slowed. The number is a dated estimate under a specified poverty definition, not a count of people who lack every physical resource. The same report describes the joint challenge of poverty, unequal opportunity and climate risk.\cite{intro-worldbank}

The WHO's 2025 report on social determinants of health equity documents how conditions of birth, work, housing and access to power, money and resources shape health. Its comparison of life expectancy across places and social groups shows why a national average cannot stand in for everyone's lived condition.\cite{intro-health} A resource model that reports only a total can miss the distribution that determines whether a person can use the resource at all.

A stock can be adequate for a district as a whole and inadequate for one clinic. A road can be passable to people with vehicles and effectively closed to people who cannot travel. A treatment can exist and remain unavailable to a patient. The relevant physical account must have enough structure to keep these distinctions visible; a total alone is often too coarse.

\section{Resources, grievance and conflict}
Control of land, fuel, water and revenue can become part of a conflict. Exclusion and inequality can deepen grievances; damaged services can make an already fragile situation worse. The World Bank identifies economic exclusion among factors that can feed fragility, alongside governance and other causes.\cite{intro-conflict} The IPCC finds that climate variability and extremes can affect conflict risk in some contexts, but its assessment warns against attributing armed conflict to climate change alone; the link is conditional and weaker than major social and political drivers.\cite{intro-ipcc-conflict}

The same care applies to the claim that monetary gain causes war. Armed conflicts have many histories and motives. Material advantage can be one incentive and control of resources one means of sustaining violence; neither statement is a universal explanation. A book about physical accounting should make the connection discussable without turning a grave human event into a slogan.

For our small district, imagine two towns drawing from one river during a drought. An account showing each withdrawal and the river's changing condition will not settle their rights or eliminate hostility. It may, however, prevent a negotiation from proceeding as though each town possessed the entire flow. Shared facts do not choose a just rule, but a just rule is harder to exercise when the facts are invisible.

\begin{intuitionbox}{Three shortages}
Not enough usable resource, no practical route to it, and no effective right to receive it are different failures. A responsible response begins by finding which failure is present.
\end{intuitionbox}

\section{From a district to a planet}
Two towns can sometimes import what they lack. Humanity as a whole cannot import a second biosphere on the schedule of a procurement contract. That does not make every natural process a fixed stock; sunlight, water cycles and regeneration are continual processes. It does make their condition and limits central to the future. The next chapter moves from the local access problem to measured planetary pressures and asks what the present evidence says about the world we are passing on.
